\documentclass[10pt]{article}
\usepackage[margin=0.91in]{geometry}
\usepackage{amsmath,amssymb,amsthm,hyperref}
\hypersetup{colorlinks=true,linkcolor=blue,citecolor=blue,urlcolor=blue}
\newtheorem{theorem}{Theorem}
\DeclareMathOperator{\diam}{diam}
\title{The integer-budget endpoint for graph diameter\protect\\ under an inverse-degree constraint}
\author{Shaitetsu}
\date{23 November 2026}
\begin{document}
\maketitle
\begin{abstract}
Under a fixed inverse-degree budget $b>1$, the extremal diameter has leading coefficient $2(\lceil b\rceil-1)$ on the scale $\log n/\log\log n$.
The companion manuscript already obtains the matching lower bound for every budget and the sharp result for noninteger budgets. The decisive complementary step is the integer-budget upper bound, obtained by charging a small endpoint at each marked crossing of the logarithmic layer-size profile. We include the whole proof and the clique-chain construction for completeness.
\end{abstract}

For a finite simple connected graph $G$ of order at least two, set
\[
 J(G)=\sum_{v\in V(G)}\frac{1}{\deg_G(v)},\qquad
 D_b(n)=\max\{\diam(G):|V(G)|=n,\ J(G)\le b\}.
\]
For fixed $b>1$ the maximum is defined for all sufficiently large $n$, since $J(K_n)=n/(n-1)$. Natural logarithms are used throughout. Erd\H{o}s, Pach and Spencer \cite{eps} studied both this diameter problem and the distinct problem of the largest \emph{mean} distance under the same constraint. Their 1988 upper bound on the diameter has leading coefficient $6b$ in their stated fixed-budget range $b\ge3$. The 2008 bound of Dankelmann, Swart and van den Berg \cite{dsv,li} gives $3b+2$. The supplied companion manuscript \cite{companion}, Theorem~1.1, gives a lower coefficient $2(\lceil b
ceil-1)$ and an upper coefficient $2\lfloor b
floor$. These agree off the integers. At an integer $b$, its bounds leave the interval $[2b-2,2b]$, and its Remark~5.1 predicts the lower endpoint to be sharp. The contribution of this note is an independent complementary upper-bound argument that closes this integer endpoint. The overlapping noninteger conclusion and lower bound are presented for completeness, not as new results relative to \cite{companion}.

\begin{theorem}\label{thm:main}
For every fixed $b>1$,
\[
 \lim_{n\to\infty}\frac{D_b(n)\log\log n}{\log n}
 =2(\lceil b\rceil-1).
\]
In particular the coefficient is $4$ when $b=3$.
\end{theorem}

\section{Layer profiles and the upper bound}
Fix a graph with $J(G)\le b$, choose one endpoint of a diametral pair, and let $s_0,\ldots,s_d$ be its breadth-first layer sizes, with $d=\diam(G)$ and $s_{-1}=s_{d+1}=0$. Every nonempty layer is represented. Edges join vertices in the same or adjacent layers, so
\begin{equation}\label{eq:budget}
 w_i:=\frac{s_i}{s_{i-1}+s_i+s_{i+1}-1},
 \qquad \sum_{i=0}^{d} w_i\le J(G)\le b .
\end{equation}
Write $x_i=\log s_i/\log n\in[0,1]$. A weak local maximum of the size profile has $w_i\ge1/3$, using the exterior zeros at the ends. Thus the number of weak local maxima is at most $3b$. For a height $h\in[0,1]$, call each maximal interval of indices with $x_i>h$ a superlevel component. Each contains a weak local maximum: choose an index attaining its maximum, and note that the adjacent indices outside the component are lower. Thus at most $6b$ edges of the index path cross any height other than a layer height. Integrating gives
\begin{equation}\label{eq:variation}
 \sum_{i=0}^{d-1}|x_{i+1}-x_i|\le6b .
\end{equation}

To extract a sharper count, consider sequences for which $d\to\infty$; bounded $d$ gives zero in the limiting normalization. Call index $i$ selected if $w_i\le (\log d)^2/d$. By \eqref{eq:budget} there are at least $d+1-bd/(\log d)^2$ selected indices. For each selected $i$ choose a neighbor of largest size. Its size is at least $R_d s_i$, where
\[
 R_d=\frac12\left(\frac{d}{(\log d)^2}-1\right);
\]
this follows by solving the bound on $w_i$ for the sum of the two neighboring layer sizes. For large $d$, $R_d>1$, so each chosen edge is identified by its unique smaller endpoint, and no chosen edges repeat. Each spans a height interval of length at least $\ell=\log R_d/\log n$.

Fix $0<\eta<1$. Discard chosen edges with height span greater than $\eta/2$, at most $12b/\eta$ edges by \eqref{eq:variation}. Inside the span of each remaining edge mark any subinterval of length $\ell$, disregarding its endpoints when counting crossings. At a marked height $h$, the selected smaller endpoint is outside \emph{all} superlevel components at $h$. Both neighbors of that endpoint have size at most $n^{\eta/2}$ times its size, since the chosen neighbor was the larger one. The smaller endpoint therefore contributes at least
\begin{equation}\label{eq:extra}
 w_i\ge\frac{1}{1+2n^{\eta/2}}\ge\frac{n^{-\eta/2}}{3}.
\end{equation}

Now remove the heights within $\eta$ below any weak local maximum; the union of these bands has measure at most $3b\eta$. At an unremoved height $h$, any superlevel component $B$ has a maximum of height at least $h+\eta$. If $S_B=\sum_{i\in B}s_i$, then $S_B\ge n^{h+\eta}$. The two exterior neighboring layers each have size at most $n^h$. Bounding each denominator in \eqref{eq:budget} on this component by $S_B+2n^h$ gives
\begin{equation}\label{eq:block}
 \sum_{i\in B}w_i\ge\frac{S_B}{S_B+2n^h}\ge1-2n^{-\eta}.
\end{equation}
Write $m=\lceil b\rceil-1$. If $b$ is not integral, then $m+1>b$ and for large $n$ the budget cannot pay for $m+1$ such components. If $b$ is integral, the bound from \eqref{eq:block} alone cannot exclude $b$ components. But at an unremoved height with a marked crossing we may charge one additional smaller endpoint using \eqref{eq:extra}, disjoint from the $b$ components. For large $n$ the resulting lower bound would be
\[
 b(1-2n^{-\eta})+\frac{n^{-\eta/2}}3>b.
\]
Hence there are at most $m$ components at any unremoved height with a marked crossing, for both integral and nonintegral $b$. There are then at most $2m$ marked edges crossing that height. If there are no marked crossings the same bound is automatic. On the removed bands we still have the crude crossing bound $6b$.

Integrating the number of marked intervals now gives
\[
 \left[d+1-\frac{bd}{(\log d)^2}-\frac{12b}{\eta}\right]
 \frac{\log R_d}{\log n}\le 2m+18b^2\eta .
\]
For fixed $\eta$, the left-hand prefactor is $d(1+o(1))$ and $\log R_d\sim\log d$. In particular $d\log d=O_b(\log n)$, choosing any fixed $\eta$. On any subsequence with normalized diameter $d\log\log n/\log n$ bounded below by a positive constant, this implies $\log d\sim\log\log n$: the lower bound on $d$ is a constant times $\log n/\log\log n$, and the upper bound is $O_b(\log n)$. Consequently the limsup of the normalized diameter is at most $2m+18b^2\eta$; then let $\eta\downarrow0$. This proves the upper bound for Theorem~\ref{thm:main}.

\section{The matching construction at all sufficiently large orders}
The companion manuscript already has this lower coefficient, including at integer budgets. We include a direct construction and an all-orders padding argument so this note can be read on its own.
For $k\ge2$ put $q=\lceil k(\log k)^2\rceil$. Take $m$ consecutive profiles
\[
 1,q,q^2,\ldots,q^{k-1},q^k,q^{k-1},\ldots,q^2,q,1,
\]
identifying the last layer of each profile with the first layer of the next one. Replace every layer by a clique of the indicated size, and add all edges between consecutive layers, with no other edges. The resulting connected graph has $2mk+1$ layers. Between distinct layers, distances equal the difference of their indices, so its diameter is $2mk$.

Every nonpeak layer has a neighboring layer at least $q$ times its size. Its contribution to $J$ is at most $1/q$, because each degree is at least the size of that neighboring layer. Each peak layer contributes at most $q^k/(q^k-1)$, by its internal clique alone. Consequently
\[
 J\le m\frac{q^k}{q^k-1}+\frac{2mk+1-m}{q}=m+o(1)<b
\]
for all sufficiently large $k$. The final inequality uses $m<b$, also when $b$ is integral. Write $N_k$ for its total order. The profile has maximum $q^k$ and $O_b(k)$ layers, so $\log N_k\sim k\log k$. Also $N_k$ increases with $k$ for large $k$, directly from the profiles and the nondecrease of $q$.

For any sufficiently large specified $n$, choose the largest $k$ with $N_k\le n$. Enlarge one peak clique by $n-N_k$ vertices and preserve its complete joins. The peak's upper bound $z/(z-1)$ only decreases as its size $z$ increases, and every nonpeak layer still has its required larger neighbor. Thus the same budget bound holds. The diameter is still $2mk$, and $N_k\le n<N_{k+1}$ implies $\log n\sim k\log k$, hence $k\sim\log n/\log\log n$. This proves the matching liminf and completes the proof.

\paragraph{Scope.} The fixed-$b$ asymptotic does not assert uniformity for a budget growing with $n$. Since mean distance is at most diameter, the upper bound also bounds the mean, but there is no claim that its coefficient is sharp for mean distance. The displayed proof is self-contained. The complementary integer upper bound was developed independently of the companion manuscript and audited. The companion manuscript treats maximum mean distance separately in its Theorem~1.2; we make no improvement claim here for that invariant. The earlier bounds cited above establish context; a comprehensive priority search and external refereeing are not claimed.

\begin{thebibliography}{9}
\bibitem{eps} P. Erd\H{o}s, J. Pach and J. Spencer, \emph{On the mean distance between points of a graph}, Congressus Numerantium \textbf{64} (1988), 121--124. \href{https://www.renyi.hu/~p_erdos/1988-23.pdf}{Author archive}.
\bibitem{dsv} P. Dankelmann, H. C. Swart and P. van den Berg, \emph{Diameter and inverse degree}, Discrete Mathematics \textbf{308} (2008), 670--673. \href{https://doi.org/10.1016/j.disc.2007.07.053}{DOI}.
\bibitem{li} X. Li and Y. Shi, \emph{On the diameter and inverse degree}, Ars Combinatoria \textbf{101} (2011), 481--487. \href{https://combinatorialpress.com/article/ars/Volume%20101/volume-101-paper-36.pdf}{Publisher PDF}.
\bibitem{companion} Anonymous, \emph{The diameter of graphs with given inverse degree: The constant in a theorem of Erd\H{o}s, Pach and Spencer}, supplied unrefereed companion manuscript, dated 24 November 2026, six pages, Theorems 1.1--1.2 and Remark 5.1 (file \texttt{Inverse\_Degree\_Diameter.pdf}).
\end{thebibliography}
\end{document}
